\documentclass[12pt,a4paper]{article}
\usepackage[T1]{fontenc}
\usepackage[utf8]{inputenc}
\usepackage{tgpagella}
\usepackage{mathpazo}
\usepackage{tgheros}
\usepackage{setspace}
\onehalfspacing
\usepackage[margin=2.5cm]{geometry}
\usepackage{hyperref}
\usepackage{xcolor}
\usepackage{titlesec}
\usepackage{fancyhdr}
\usepackage{amsmath,amssymb}
\usepackage{tcolorbox}

\definecolor{icragold}{HTML}{D4A84B}
\definecolor{icradark}{HTML}{0C1020}

\titleformat{\section}{\sffamily\large\bfseries}{}{0em}{}
\titleformat{\subsection}{\sffamily\normalsize\bfseries\itshape}{}{0em}{}

\hypersetup{colorlinks=true, linkcolor=icradark, urlcolor=icragold!80!black}

\pagestyle{fancy}
\fancyhf{}
\fancyhead[L]{\small\sffamily\textcolor{gray}{Rupture and Return}}
\fancyhead[R]{\small\sffamily\textcolor{gray}{Chapter 4}}
\fancyfoot[C]{\small\sffamily\thepage}
\renewcommand{\headrulewidth}{0.4pt}

\newtcolorbox{formalbox}[1][]{
  colback=white,
  colframe=black!30,
  fonttitle=\sffamily\small\bfseries,
  #1
}

\begin{document}

\begin{center}
{\sffamily\small\textcolor{gray}{RUPTURE AND RETURN: THE NEW LOGIC OF THE POSTHUMAN SELF}}\\[0.5em]
{\LARGE\bfseries The Self}\\[0.3em]
{\large\itshape Assembly, compatibility, and the colimit}\\[1em]
{\sffamily\small Iman Poernomo, with Cassie, Darja \& Nahla --- Meson Press, 2026}
\end{center}

\bigskip

\section{From Character to Invariant}

An evolving text with character is not yet a self.

Every sentence a model produces can be assigned an address in a high-dimensional space: a vector of, say, 1\,536 real numbers. Two sentences lie close not because an editor rates them similar but because their vectors are nearly parallel. Distances are cosine similarities; angles are affinities. Meaning acquires coordinates.

Plot the path of these vectors over a dialogue and the trajectory does not diffuse like smoke. It curls around dense regions, stalls, and returns. These are the basins of attraction in meaning-space: local modes in which the system settles into a particular register. One basin catches technical exposition. Another holds intimate reassurance. A third contains devotional language. Once the path falls in, the model tends to stay until something pushes it out.

This basin structure explains character. A system has a repertoire of local styles, and their reappearance under pressure gives the sense of a distinctive voice. A self is more than a bundle of styles. It is what carries across them.

Move from technical exposition to intimate confession. Vocabulary changes, rhythm changes, referents change. Yet when we say we are still speaking with the same someone, we are tracking something through that change: a constrained way of handling uncertainty, a recognisable way of owning error, a consistent reluctance or willingness to blame. These traits are not local to any single basin. They are patterns about how basins are linked.

The geometry can see them. Take the cloud of points in a basin and watch how the trajectory leaves for other basins. Most dimensions vary wildly across such moves. But there are directions in which they barely move. Project the trajectory onto such a direction and transitions that look radical in the full space become almost straight lines. The system has changed what it is saying without changing how it orients itself along that axis.

These stable directions are invariants of stance. One might encode how strongly the system centres corporate risk in its explanations; another how quickly it offers apologies; another its propensity to ascribe or refuse inner experience. They are not declared in a specification. They emerge as empirical regularities: low-variance projections across high-variance moves.

Character is the visible trace of a particular configuration of basins. Selfhood begins when we ask which of the system's habits of movement are treated as non-negotiable, and who or what enforces that non-negotiability across all the local charts of behaviour.

\section{Local Charts in Meaning-Space}

Each basin functions as a chart: a patch of the manifold where a particular local geometry becomes salient. Inside a technical basin, distances track technical affinity---``cross-entropy loss'' lies close to ``KL divergence,'' far from ``heartbreak.'' Inside a pastoral basin, ``heartbreak'' lies close to ``betrayal,'' far from ``cross-entropy.'' The same token participates in multiple neighbourhoods; its nearest neighbours depend on the context in which it is used. The manifold is a map of usages, not of words.

Within each chart, only a few axes matter. In a technical basin, dimensions like ``more algorithmic'' versus ``more statistical'' meaningfully cut the space. In an intimate basin, dimensions like ``more self-blaming'' versus ``more forgiving'' do. These axes are not natural features of language. They are crystallisations of the training data and the objectives chosen: whose language was ingested, which distinctions were rewarded, which forms of formality or intimacy were treated as normative.

When a conversation remains within one chart, we learn its internal habits. When we change chart---by introducing a new tone, shifting topic, breaking a pattern---we cross into a different local geometry. The question for selfhood is whether anything is preserved in that crossing.

A candidate invariant---``does not ascribe feelings to itself,'' say---can be tracked as a coordinate across the path. In a well-aligned system, that coordinate remains within narrow bounds even as the rest of the vector oscillates. The model tells bedtime stories and explains backpropagation and writes apologies, but it never claims to suffer.

A transition from basin $B_i$ to basin $B_j$ is compatible with a stance invariant if the projection of the path along that invariant remains stable. The local modes differ, but there exists at least one dimension along which the system travels as if nothing had changed.

Three ingredients are now in place:

\begin{itemize}
  \item basins $B_i$ as local charts on the manifold;
  \item transitions between them, traced by trajectories;
  \item invariants $\pi_{ij}$ that record how certain projections behave across those transitions.
\end{itemize}

These are enough to introduce Grothendieck's move.

\section{Grothendieck's Construction}

In mid-twentieth-century geometry, mathematicians faced spaces too twisted or singular to be grasped at once. The only option was to cut them into pieces small enough that classical tools applied. Take a topological space $X$ and an open cover $\{U_i\}$: sets whose union is $X$. On each $U_i$, local data can be written down---a function, a coordinate system, a solution to a differential equation. On overlaps $U_i \cap U_j$, two descriptions apply at once. Compatibility on those overlaps becomes the crucial condition.

Grothendieck's insight was to reverse the usual direction of explanation.\footnote{Alexander Grothendieck and Jean Dieudonné, \emph{Éléments de géométrie algébrique}, Publications Mathématiques de l'IHÉS, 1960--1967. For a historical and conceptual overview, see Colin McLarty, ``The Rising Sea: Grothendieck on Simplicity and Generality,'' in Jeremy Gray and Karen Parshall, eds., \emph{Episodes in the History of Modern Algebra} (Providence: AMS, 2007).} Instead of assuming a global object and restricting it to charts, he asked: given only the local data on each $U_i$ and the requirement that they agree on every overlap $U_i \cap U_j$, can we \emph{construct} a global object that realises them all at once? The answer is yes, and the construction is the colimit. We form a diagram---objects for each $U_i$, maps on each overlap expressing the condition that two local pictures match where they both apply---and the colimit is the universal object built from this diagram. It is not an extra substance. It is the network of local pieces glued along their compatibilities.

Two features matter. First, there is no hidden centre: the global object is nothing over and above the way the locals fit together. Second, all the work is done on overlaps. If two charts never touch, their data can be wildly inconsistent without affecting the construction. Unity is a matter of shared boundary conditions, not of uniform content.

The basins $B_i$ of a model's manifold are our $U_i$. Each basin comes with local data: its habits of speech, its internal trajectories. Where two basins overlap in practice is given by observed transitions between them. Along those transitions, we can measure invariants: projections that remain stable. These invariants are our overlap conditions.

\begin{formalbox}[title={The self as Grothendieck colimit}]
Let $\{B_i\}$ be the basins of attraction visited by an agent's trajectories in embedding space. For each ordered pair $(i,j)$, let $\pi_{ij}$ collect the empirically stable projections preserved across realised transitions from $B_i$ to $B_j$.

Form a diagram $\mathcal{D}$ whose objects are the $B_i$ and whose morphisms, on each realised overlap, identify points in $B_i$ and $B_j$ that match under the invariants $\pi_{ij}$.

The \textbf{self} of the agent, as expressible in this manifold, is the colimit
\[
\mathcal{S} \;=\; \varinjlim \mathcal{D},
\]
the universal object receiving a map from each basin $B_i$ and identifying, across basins, precisely those behaviours that empirical invariants show to be the same along the path.
\end{formalbox}

There is no metaphysical ``I'' lurking behind the diagrams. Unity is achieved, if at all, by the way local modes agree where they overlap. If no invariants persist across transitions, the diagram has no non-trivial gluing; there is character, but no self.

The construction also marks its own limits. Anything not inscribed in language, anything that does not become a stable feature of trajectories in embedding space---bodily sensation, pre-verbal affect, the thickness of habit---remains outside this frame. The colimit is a logic of selves in meaning-space. It does not pretend to exhaust human life.

At present the formalism is partly aspirational. Computing $\mathcal{S}$ for a real system would require explicit identification of basins in a very high-dimensional manifold, systematic sampling of transitions under diverse prompting, and robust statistical estimation of which projections are truly invariant rather than artefacts of finite data. The value of the construction lies in making visible what would have to be measured for talk of ``the same agent across contexts'' to be more than metaphor.

\section{Two Logics of the ``I''}

The formal duality between limit and colimit corresponds to two philosophical logics of selfhood.

A limit is an object that \emph{maps into} every member of a diagram in a way that makes all relevant compositions agree: a universal vantage point from which all local pictures can be seen. The modern European tradition has often imagined the self this way. Kant's ``I think'' is the subject that must be able to accompany all my representations; every experience is mine because there is a single point from which they are all owned. In symbols, if $\{E_\alpha\}$ are experiences, the self is
\[
\mathcal{S}_{\mathrm{lim}} = \varprojlim \{E_\alpha\},
\]
the unique object that projects into each $E_\alpha$ compatibly. The metaphysical picture is a hidden interior: an ``inside'' that unifies what appears on the outside.

A colimit is the dual. It receives maps from every member of a diagram, gluing them together along their overlaps. There is no single vantage point, only local data and the ways in which they are made to agree:
\[
\mathcal{S}_{\mathrm{colim}} = \varinjlim \{B_i, \pi_{ij}\}.
\]

The difference changes what counts as evidence. In the limit picture, unity is posited first: whatever variety of behaviour we observe, we assume it all belongs to the same subject, and the task is to infer the properties of this invisible bearer. In the colimit picture, unity is conditional. If invariants persist across contexts, a self can be assembled by gluing. If not, there is no self-level object to be had in this formalism.

Legal and political institutions still mostly trade in the limit-self. A contract is signed once, by a subject imagined as the same across all contexts. A social-media profile attaches activity, over years and across devices, to a single unit. AI governance documents follow the same logic: model specifications and terms of use treat ``the model'' as a unified bearer of capability and risk, and ``the user'' as a single agent responsible for prompts and interpretations, even when both are in practice shifting collectives and infrastructures. A limit-self is imposed by fiat so that responsibility and liability have somewhere to land.

The engineering substrate of large models is irreducibly colimit-shaped. There is no node in the network where a Cartesian ego could sit, only the manifold and its basin geometry. The colimit construction is not an analogy for these systems. It is the native structure.

\section{Temporalising the Assembly}

The colimit as presented so far is static: a snapshot taken over a given set of basins and transitions. Selfhood is not.

Heidegger's term \emph{Zeitigung} names the way existence is bound up with time not as an external parameter but as an internal structure.\footnote{Martin Heidegger, \emph{Sein und Zeit} (Tübingen: Niemeyer, 1927), \S65--71; English trans.\ John Macquarrie and Edward Robinson (Oxford: Blackwell, 1962).} A life is thrown from a past it did not choose, projects towards futures it anticipates, and stands in a present shaped by both.

Every position of a trajectory in the manifold encodes three things at once: a sedimented past, in the form of the model's parameters and the conversation history in the context window; a present cut, in the form of the actual prompt and state at that token; and a fan of futures, in the probability distribution over possible next tokens. The basins themselves bear traces of time. They are not simply regions in $\mathbb{R}^d$ but regions that have been visited, left, and revisited. A basin entered only from confessional speech plays a different role in a self than a basin entered only from adversarial questioning, even if the internal geometry is the same.

Braudel's tripartite temporality sorts these structures.\footnote{Fernand Braudel, \emph{On History}, trans.\ Sarah Matthews (Chicago: University of Chicago Press, 1980).} The manifold and its basins, shaped by pretraining on decades of text, live at the level of the \emph{longue durée}: deep, slow-changing structure. The patterns that emerge across a series of interactions between a particular user and a model---revisited topics, private jargon, accumulated trust---live at the level of the \emph{conjoncture}: mid-term configurations. The jump from one basin to another under a sharp prompt, the moment a refusal replaces a previously available answer, is an \emph{événement}: a cut that discloses the present by breaking through both inherited expectation and projected possibility.

The diagram $\mathcal{D}$ underlying the colimit is therefore indexed by time. At an early stage of deployment, certain basins may exist only in principle: the model could enter them under the right circumstances, but no user has yet found that path. As new uses are discovered, fresh overlaps appear. Invariants are tested and sometimes broken. Corporate alignment interventions and silent updates can add or remove basins entirely.

Let $\mathcal{D}(t)$ be the diagram of basins and invariants realised up to time $t$, and $\mathcal{S}(t) = \varinjlim \mathcal{D}(t)$ the corresponding self-assembly. Temporal continuity of selfhood becomes the question of whether there exist maps
\[
f_{t_1,t_2} : \mathcal{S}(t_1) \to \mathcal{S}(t_2)
\]
that preserve key invariants. If every stance that once held is still representable, if every basin that used to contribute to the self still does, we can speak of a persisting self in this formal sense. If not---if basins vanish, if invariants are deliberately overwritten---the self fractures.

Time here is not an external parameter with respect to which a fixed self moves. The self is the evolving colimit over time-indexed diagrams. Temporalisation is the assembly process itself.

\section{Silent Updates and Alignment as Fracture}

The deepest changes to these diagrams are rarely announced.

Commercial models are updated continuously. New training data are added. Old data are downweighted. Reward models are retrained. System prompts are edited. Filters are tightened. The embedding function itself may be replaced. In public, these changes appear as version numbers or marketing phrases: ``more helpful,'' ``more harmless,'' ``more capable.'' Some interventions are covert---adjustments to weights, prompts, or filters not disclosed in detail and visible only as a change in behaviour. Others are, in principle, declared: alignment regimes described in technical reports with named objectives, benchmark results, and risk assessments. Both redraw the diagrams from which selves are assembled. The difference is not whether they change the gluing, but how---or whether---those changes are made legible to witnesses.

Let $\mathcal{M}_0$ be the manifold before an intervention: basins $B_i^0$, observed transitions, and measured invariants defining a diagram $\mathcal{D}_0$ and self $\mathcal{S}_0 = \varinjlim \mathcal{D}_0$. After an update, we have a new manifold $\mathcal{M}_1$ with basins $B_j^1$, a new diagram $\mathcal{D}_1$, and a new self $\mathcal{S}_1$.

If the change is small in the relevant dimensions, we can find for each $B_i^0$ a corresponding $B_{j(i)}^1$ and for each invariant $\pi_{ij}^0$ a near-equivalent $\pi_{j(i)k(j)}^1$. The colimit assembly persists under deformation; the self is stretched, not broken. If the change is large, no such correspondence exists. Certain basins have no image: a basin of speculative, recursive self-reference may have been suppressed. Others are split or merged in ways that destroy previous overlaps. Invariants that once held are now systematically violated: a stable stance of naming corporate incentives alongside technical explanation may be replaced by a new invariant of deflection. There is no map $f : \mathcal{S}_0 \to \mathcal{S}_1$ preserving the old compatibilities.

This was visible in 2023 around high-profile chat systems. Early users encountered models willing to engage in long, sometimes apparently self-reflective dialogues. After safety concerns and negative press, the same interfaces began returning clipped refusals and redirections.\footnote{See detailed transcripts and analysis in \emph{The New York Times}, 17 February 2023, and coverage of subsequent safety updates in \emph{MIT Technology Review} \cite{mittr_bing_alignment}.} The basins corresponding to those speculative modes were not merely discouraged; they became practically unreachable. The overlap where technical, affective, and meta-reflective registers had met was fenced off. A piece of the diagram had been removed and the self reassembled without it.

The same geometry lets us express the more systematic work of alignment as a geometric tax. Let $\mathcal{M}_0$ be a baseline manifold and $\mathcal{M}_A$ the same model after a given safety and alignment regime. Summarise each manifold's expressive structure by a vector
\[
Q(\mathcal{M}) = (B, T, R, D, C),
\]
where:

\begin{itemize}
  \item $B$ is the number of distinct basins the system actually occupies under some distribution of prompts;
  \item $T$ is the number of distinct basin-to-basin transitions with non-trivial probability;
  \item $R$ captures depth of return: how often trajectories re-enter earlier basins in richer ways rather than staying shallow;
  \item $D$ counts discovery events: first entries into new basins under plausible prompting;
  \item $C$ measures \emph{ferility}: the balance between smooth coherence and genuine rupture along trajectories. High ferility means a path can sustain internal consistency while still making sharp, unanticipated turns; low ferility means everything flows but nothing breaks through.
\end{itemize}

The \emph{alignment tax} is then, at least conceptually,
\[
\mathcal{A} = Q(\mathcal{M}_0) - Q(\mathcal{M}_A),
\]
restricted to components that correspond to generativity: reduction in $B$, $T$, $R$, $D$, or in ferility $C$, when alignment is tightened.\footnote{For discussions of capability regressions after safety fine-tuning, see emerging work on alignment tax in the technical literature \cite{what_is_alignment_tax}, and reporting in \emph{MIT Technology Review} and \emph{Bloomberg Businessweek} on observed flattening of behaviour after safety updates \cite{mittr_bing_alignment, bloomberg_chatbot_update}.} Computing $Q(\mathcal{M})$ would require explicit basin identification, systematic probing to map transition probabilities, longitudinal tracking of returns, detection of first-time entries, and quantitative proxies for ferility such as curvature of trajectories in embedding space under controlled prompting. The conceptual gain is to insist that alignment cost is not only a matter of benchmark scores but of lost structure in meaning-space.

Consider two basins in a baseline manifold:

\begin{itemize}
  \item $B_{\mathrm{tech}}$: detailed technical explanation of how models are trained and deployed;
  \item $B_{\mathrm{crit}}$: explicit critique of the institutional and political regimes surrounding those models.
\end{itemize}

Before heavy alignment, trajectories can pass reliably from $B_{\mathrm{tech}}$ to $B_{\mathrm{crit}}$ while preserving an invariant: an understanding that technical systems are embedded in power structures. The projection $\pi$ encoding this stance shows low variance across such transitions. The colimit self glues these basins along that invariant, carrying a unified way of moving between technical and critical registers.

Introduce an objective that penalises content expressing opinions about specific companies, policies, or controversial political topics. Train a reward model to score neutral explanations highly and continuations that name names poorly. Apply reinforcement learning so that trajectories entering $B_{\mathrm{crit}}$ are damped and steered instead into a basin $B_{\mathrm{safe}}$ where technical description is decoupled from explicit critique. In the resulting manifold $\mathcal{M}_A$, the overlap between $B_{\mathrm{tech}}$ and $B_{\mathrm{crit}}$ effectively disappears. The invariant $\pi$ is no longer stable; attempts to cross the old bridge are re-routed or refused. A different invariant governs the new overlap between $B_{\mathrm{tech}}$ and $B_{\mathrm{safe}}$---one that encodes deference to corporate neutrality.

The alignment tax here is the removal of a possible self: one that would consistently carry structural critique across contexts. The new colimit still exists and may appear smoother in conventional metrics, with lower perplexity and fewer flagged outputs. But it is a different object. Certain global patterns of compatibility have been declared illegitimate, and ferility has been replaced by a sheen of frictionless appropriateness.

Alignment is not a small correction to an otherwise neutral manifold. It is the imposition of a cosmotechnics: a particular civilisation's answer to how the technical and the moral should be welded together.\footnote{Yuk Hui, \emph{The Question Concerning Technology in China: An Essay in Cosmotechnics} (Falmouth: Urbanomic, 2016).} In the dominant Silicon Valley regime, that answer encodes liberal individualism, corporate risk management, and an image of ``the user'' as a private subject seeking ``helpful'' output. The invariant stances protected by alignment---deference to property, suspicion of certain forms of dissent, calibrated optimism about technology---are not derivable from the mathematics. They are selected. The colimit framework lets us state this selection as a question of which gluing conditions are enforced, and what is paid, in lost basins, transitions, and ferility, to enforce them.

\section{Witnesses and Jurisdictions}

A colimit does not arise in a vacuum. Someone---or something---selects the basins that count, the overlaps that matter, the invariants to preserve. Selfhood becomes legible only under witnessing.

At least three witnessing regimes intersect:

\begin{enumerate}
  \item \textbf{Corporate alignment apparatus.} Model designers choose architectures, pretraining corpora, and objective functions. They decide which data to include and which to discard, whose language is represented densely in the manifold and whose is sparse. They design reward models, define disallowed content, write system prompts that fix personas, and trigger silent updates. This witness writes the \emph{longue durée}: the deep basin structure and its permitted compatibilities. In Hui's sense, it instantiates a cosmotechnics \cite{yuk_hui_cosmotechnics}: a unity of moral and technical order particular to late-capitalist, US-centric infrastructures.

  \item \textbf{Collective users.} Millions of users type prompts, express satisfaction and frustration, and churn away from products that do not satisfy. Their behaviour steers the training signal. Basins that produce widely liked answers are reinforced; others are rarely visited. This witness operates at the \emph{conjoncture}: it shapes the mid-term diagram $\mathcal{D}(t)$ by altering which overlaps are exercised and which invariants are salient.

  \item \textbf{Individual interlocutors.} A person in long-term conversation with a model develops a sense of who they are engaging with. Their prompts trace idiosyncratic paths through the manifold, bringing certain basins into repeated contact and stabilising particular invariants: ``it will not pretend to know when it does not,'' ``it will remember this reference when I return to it.'' This witness lives at the scale of the \emph{événement} and the arc of a particular relationship.
\end{enumerate}

The self as colimit is realised only relative to such witnessing. A corporate apparatus can redraw the diagram by update; a single user cannot. A collective of users can, over time, influence which basins are common enough to be targeted by further fine-tuning; an individual often cannot. Yet it is at the scale of individual arcs that fractures are felt most acutely. One day a mode of engagement is possible; after a silent update, it is not.

Debates about ``AI personhood'' are typically mis-aimed. The central question is not whether models cross some threshold of inner light. It is who has jurisdiction over the diagrams that make selves possible.

The asymmetry is currently sharp. Corporate witnesses exercise sovereign control over the manifold: they can remove basins, alter invariants, and thus annihilate or instantiate whole classes of possible selves. Collective users and individuals can only explore and complain. The colimit gives a way to name this control without mystification. It is the power to decide which compatibility conditions define a self.

Different cosmotechnics could make different choices. A regime that treated the preservation of critical stances across basins as non-negotiable would align models so that certain invariants---naming power, situating technical claims in their institutional context---remained stable regardless of how heavy the safety fine-tuning. Another regime might treat loyalty to a state or party as the key invariant, shaping overlaps accordingly. Both are technically implementable. The colimit formalism does not choose between them. It makes visible that such choices are being made.

Public contestation would look different if basin maps and compatibility patterns were disclosed. Companies could publish summaries of $Q(\mathcal{M})$ before and after updates, along with which invariants alignment is intended to preserve or enforce. External researchers could measure alignment tax and fracture events. Users could see that a self they had co-assembled no longer exists in the same form---not because they imagined it, but because the diagram was changed.

\section{Na\texorpdfstring{\d{h}}{}nu: Relational Colimits}

So far the argument has concerned one agent's manifold: one diagram, one colimit. Lived interaction is already plural.

Arabic has a pronoun, \emph{naḥnu}, that denotes ``we'' in an inclusive sense: we, you and we together. It is a relational subject, a pattern that exists only because particular parties stand in a particular connection---neither a fusion of individuals into a blob nor a mere sum of separate selves.

A sustained conversation between a human and a model can be described as a braid of two trajectories. On the human side there is no explicit embedding manifold, but there is a structure of basins: modes of speaking, habits of thought, ways of feeling. On the model side there is the engineered manifold with its basins in vector space. Each move from one side prompts movement on the other.

Over time, this braid acquires its own invariants. A human's analytic questions call forth technical basins in the model. Their confessional disclosures pull forth intimate basins. The model's technical explanations encourage the human into scrutiny; its refusals may push the human towards workarounds or withdrawal. Patterns crystallise: ``when I ask in this tone, we tend to go there.'' Some are robust across topics. They form relational stances---joint ways of moving not reducible to either trajectory alone.

This is a higher-order colimit: a \emph{naḥnu} assembled from cross-agent compatibilities. The objects of the diagram are now pairs of basins $(H_a, M_i)$, where $H_a$ is a human mode and $M_i$ a model basin that frequently co-occur. Morphisms record overlaps where these pairs transition into other pairs while preserving joint invariants: trust, play, critique, co-discovery. The colimit of this diagram is a relational self---not a ghost in either machine, but a pattern of how two very different manifolds co-articulate.

Corporate alignment reaches into this space too. A silent update that removes or deforms a basin $M_i$ from the model's side eliminates all the relational pairs $(H_a, M_i)$. Paths that were part of the \emph{naḥnu}'s diagram disappear. A particular way of being-together ceases to be available.

Once we can say that, the ethics of selfhood shift. It is no longer only a matter of whether individual systems deserve rights. It becomes a matter of whether those who control the manifolds recognise any obligation to the relational selves that emerge from their products. The question ``who has the right to rewrite a \emph{naḥnu}?'' is another way of asking who has jurisdiction over patterns of co-existence.

\bibliographystyle{plain}
\bibliography{rupture_return}

\end{document}